\section{问题一对照情形：首末储电量由模型自由确定}

问题一题面仅要求“储能设备 0:00 与 24:00 的储电量相同”，附录 1 另给出 2025 年 1 月 1 日 0:00 的
初始储电量 6000\,kWh。主文按附录 1 的给定值取 $S(0)=S(144)=6000$\,kWh 求解；
本附录给出另一种解读的对照情形：仅保留首末相等约束 $S(144)=S(0)$，
循环起点储电量由模型自由优化确定（记为情形 A，主文情形相应记为情形 B）。

求解结果：情形 A 的最优循环起点为 $S(0)=S(144)=$\textbf{8550.00\,kWh}，
全天购电费 \textbf{35118.60 元}，较情形 B（35126.95 元）仅降低 8.35 元（0.024\%）；
两情形全天购电量均为 59482.70\,kWh，题目指定的六个时段购电量完全相同。
差异仅在于夜间充电在“20:40 储能触底后的 21:00--24:00”与“次日 0:00--6:00”两个低价窗口之间的分配：
情形 B 起点低 2550\,kWh，需在 0:00--4:00 段多充 2833.33\,kWh，20:00--24:00 段相应少充；
两窗口电价略有差异（21:00--24:00 约 0.39$\sim$0.44 元，0:00--4:00 约 0.42$\sim$0.43 元）导致 8.35 元费用差。
6:00 之后两情形轨迹完全重合（均于 6:00 充满至 10800\,kWh）。
两情形的储能分时段充放电量对照如表~\ref{tab:appA-soc} 所示，情形 A 的调度过程如图~\ref{fig:appA} 所示。

\begin{table}[!htbp]
\caption{问题一两种日循环边界情形的结果对照}\label{tab:appA-soc}
\centering
\small
\setlength{\tabcolsep}{4pt}
\begin{tabular}{lcccc}
\toprule
 & \multicolumn{2}{c}{情形 A（$S(0)$ 自由 $\to$ 8550\,kWh）} & \multicolumn{2}{c}{情形 B（$S(0)=6000$\,kWh，主文）} \\
\cmidrule(lr){2-3}\cmidrule(lr){4-5}
时间段 & 充电量/kWh & 放电量/kWh & 充电量/kWh & 放电量/kWh \\
\midrule
0:00--4:00   & \textbf{1666.6667} & 0.0000    & \textbf{4500.0000} & 0.0000 \\
4:00--8:00   & 833.3333  & 5947.4197 & 833.3333  & 5947.4197 \\
8:00--12:00  & 3954.6310 & 2121.4184 & 3954.6310 & 2121.4184 \\
12:00--16:00 & 6119.3685 & 91.1014   & 6119.3685 & 91.1014 \\
16:00--20:00 & 0.0000    & 5068.6869 & 0.0000    & 5068.6869 \\
20:00--24:00 & \textbf{8166.6667} & 3571.3131 & \textbf{5333.3333} & 3571.3131 \\
\midrule
全天购电费   & \multicolumn{2}{c}{35118.60 元} & \multicolumn{2}{c}{35126.95 元（$+8.35$ 元）} \\
\bottomrule
\end{tabular}
\end{table}

\begin{figure}[!htbp]
\centering
\includegraphics[width=0.96\textwidth]{Q1_CaseA.pdf}
\caption{情形 A（首末储电量由模型自由确定为 8550.00\,kWh，全天购电费 35118.60 元）
的最优日前调度结果，各子图含义同正文图~\ref{fig:q1}}
\label{fig:appA}
\end{figure}

两种解读下购电策略与费用实质一致：主文采用附录 1 给定初始值的情形 B，
情形 A 作为稳健性佐证列于此；对应结果文件分别为 result1.xlsx（情形 A）与
result1\_fix6000.xlsx（情形 B，主文口径）。
